\documentclass[11pt, a4paper, oneside]{article}
\usepackage[ngerman]{babel}
\usepackage{worksheet}

\begin{document}
	\author{L. Bung}
	\title{Abstrakte Klassen \hspace{10cm} und Methoden}
	\subject{SAE}
	\class{E2FI}
	\maketitle
	
	\begin{hint}{Abstrakte Klassen und Methoden}
		Als typisches Beispiel für eine Vererbungsrelation hatten wir schon die Klassen \emph{Auto} und \emph{Fahrrad} mit der gemeinsamen Superklasse \emph{Fahrzeug} gesehen.
		
		Wir hatten dies als Beispiel dafür genommen, dass gemeinsame Attribute (z. B. die Farbe des Fahrzeugs) oder Methoden (z. B. \texttt{fahren()}) durch die Superklasse abstrahiert werden können.
		
		In der Realität existieren allerdings nur die Subklassen -- Autos, Fahrräder usw.
		Die Klasse \emph{Fahrzeug} ist daher eine \textbf{abstrakte Klasse} und \texttt{fahren()} eine \textbf{abstrakte Methode}.
		Es können keine Objekte der Klasse \emph{Fahrzeug} instanziiert werden und die Methode \texttt{fahren()} kann (auf der Superklasse) nicht aufgerufen werden.
	\end{hint}
	
	
	\task{Umsetzung in UML}
	
	\begin{enumerate}[label=\alph*)]
		\item Recherchieren Sie, wie man abstrakte Klassen und Methoden in UML-Diagrammen kennzeichnet.
		
		\lines[3cm]
		
		\item Entwerfen Sie ein UML-Diagramm zur folgenden Situation:
			\begin{itemize}
				\item Benachrichtigungen können entweder per E-Mail oder per SMS versendet werden.
				\item Eine Nachricht hat immer einen Empfänger.
				\item E-Mails werden von einer E-Mail-Adresse gesendet.
				\item SMS werden von einer Telefonnummer gesendet.
			\end{itemize}
			
			\boxarea[5cm]
	\end{enumerate}
	
	\task{Umsetzung in Python}
	
	\begin{enumerate}[label=\alph*)]
		\item Finden Sie heraus, wie man abstrakte Klassen und Methoden in Python umsetzen kann.
		
		\lines[2cm]
		
		\item Implementieren Sie das UML-Diagramm aus Aufgabe 1 in Python.
		
		\item Was passiert, wenn keine Subklasse die abstrakte Methode der Superklasse implementiert?
		
		\lines[2cm]
	\end{enumerate}
	
	\bonustask{Eigene Beispiele}
	
	\begin{itemize}
		\item Finden Sie weitere Beispiele für Situationen, in denen die Verwendung abstrakter Klassen und Methoden sinnvoll wäre.
		\item Erstellen Sie UML-Diagramme für die Situationen.
		\item Setzen Sie die Diagramme in Python-Code um.
	\end{itemize}
\end{document}